%; whizzy chapter -dvi
% -initex iniptex -latex platex -format platex -bibtex jbibtex -fmt fmt
% 以上 whizzytex を使用する場合の設定。
 
%     Tokyo Debian Meeting resources
%     Copyright (C) 2012 Junichi Uekawa
%     Copyright (C) 2011, 2015 Nobuhiro Iwamatsu

%     This program is free software; you can redistribute it and/or modify
%     it under the terms of the GNU General Public License as published by
%     the Free Software Foundation; either version 2 of the License, or
%     (at your option) any later version.

%     This program is distributed in the hope that it will be useful,
%     but WITHOUT ANY WARRANTY; without even the implied warranty of
%     MERCHANTABILITY or FITNESS FOR A PARTICULAR PURPOSE.  See the
%     GNU General Public License for more details.

%     You should have received a copy of the GNU General Public License
%     along with this program; if not, write to the Free Software
%     Foundation, Inc., 51 Franklin St, Fifth Floor, Boston, MA  02110-1301 USA

%  preview (shell-command (concat "evince " (replace-regexp-in-string "tex$" "pdf"(buffer-file-name)) "&"))

%%ここからヘッダ開始。

\documentclass[mingoth,a4paper]{jsarticle}
\usepackage{monthlyreport}
% 日付を定義する、毎月変わります。
\newcommand{\debmtgyear}{2016}
\newcommand{\debmtgmonth}{5}
\newcommand{\debmtgdate}{21}
% started from zero:
% (let ((year 2013) (month 7)) (+ (* (- year 2005) 12) month -1))
\newcommand{\debmtgnumber}{139}

% tikz picture の為のマクロ設定
\usepackage[dvipdfmx]{graphicx}
\usepackage{tikz}

\begin{document}

\begin{titlepage}
\thispagestyle{empty}
% タイトルページ:編集必要な部分は最初のマクロに飛ばすこと

\vspace*{-2cm}
第\debmtgnumber{}回 東京エリア Debian 勉強会資料\\
\hspace*{-2cm}
\includegraphics{image2012-natsu/dotdeb.pdf}\\
\hfill{}\debmtgyear{}年\debmtgmonth{}月\debmtgdate{}日

% ここはアップデートすること
% 全角文字にしないとフォントのサイズが合わないので注意
\rotatebox{10}{\fontsize{30}{30} {\gt 特集 ：Buffalo Linkstation 向け Debian Installer}}\\

\vspace*{-2cm}
\hfill{}\includegraphics[height=6cm]{image200502/openlogo-nd.eps}
\end{titlepage}

\newpage

\begin{minipage}[b]{0.2\hsize}
 \definecolor{titleback}{gray}{0.9}
 \colorbox{titleback}{\rotatebox{90}{\fontsize{80}{80} {\gt デビアン勉強会} }}
\end{minipage}
\begin{minipage}[b]{0.8\hsize}
\hrule
\vspace{2mm}
\hrule
\begin{multicols}{2}
\tableofcontents
\end{multicols}
\vspace{2mm}
\hrule
\end{minipage}

\dancersection{事前課題}{杉本 典充}

今回の事前課題は以下です:
\begin{enumerate}
\item PCまたはサーバ以外にDebianをインストールしたことがありますか。
\item 先の質問で「あります」とお答えした方にお聞きします。その機械をお教えください。
\item Hack Timeは何をしますか。
\end{enumerate}
この課題に対して提出いただいた内容は以下です。
\begin{multicols}{2}
{\small
\begin{prework}{ koedoyoshida }
  \begin{enumerate}
  \item $B$"$j$^$9(B
  \item raspberry-pi 1B/2/3
  \item DebianJP miniconf2016JP/PyCon JP$B=`Hw(B
  \end{enumerate}
\end{prework}

\begin{prework}{ kenhys }
  \begin{enumerate}
  \item $B$"$j$^$9(B
  \item ($B2sEz$J$7(B)
  \item ITP$B$N:n6H$r:F3+$9$k(B
  \end{enumerate}
\end{prework}

\begin{prework}{ tyamada22 }
  \begin{enumerate}
  \item $B$"$j$^$;$s(B
  \item ($B2sEz$J$7(B)
  \item DDTSS $B$^$?$O(B po-templates$B$J$I$NK]Lu(B
  \end{enumerate}
\end{prework}

\begin{prework}{ dictoss }
  \begin{enumerate}
  \item $B$"$j$^$9(B
  \item $B9uH"(BHG$B!"%i%:%Y%j!<%Q%$(B2
  \item kfeebsd$B$H(Btilegx$B$N>pJs<}=8(B
  \end{enumerate}
\end{prework}

\begin{prework}{ Charles Plessy }
  \begin{enumerate}
  \item $B$"$j$^$9(B
  \item https://wiki.debian.org/Cloud/AmazonEC2DebianInstaller
  \item GPG$B%-!<(B \& mime-support $B%Q%C%1!<%8(B
  \end{enumerate}
\end{prework}

\begin{prework}{ Roger Shimizu }
  \begin{enumerate}
  \item $B$"$j$^$9(B
  \item ARM board$B$d!"(Bvirtualbox VM$B$J$I(B
  \item debian installer $B$K(B GNU/screen $BBP1~$K4X$9$k:n6H(B
  \end{enumerate}
\end{prework}

\begin{prework}{ takaswie }
  \begin{enumerate}
  \item $B$"$j$^$;$s(B
  \item ($B2sEz$J$7(B)
  \item ALSA$B%f!<%6!<%i%s%I%i%$%V%i%j$N%Q%C%A$r=q$/M=Dj$G$9!#(B
  \end{enumerate}
\end{prework}

\begin{prework}{ henrich }
  \begin{enumerate}
  \item $B$"$j$^$;$s(B
  \item ($B2sEz$J$7(B)
  \item administrators handbook$B$N::FI(B
  \end{enumerate}
\end{prework}

\begin{prework}{ $B$D$k(B }
  \begin{enumerate}
  \item $B$"$j$^$;$s(B
  \item ($B2sEz$J$7(B)
  \item Debian $B?7%a%s%F%J!<%,%$%I$rFI$`(B
  \end{enumerate}
\end{prework}

}
\end{multicols}

\dancersection{Debian Trivia Quiz}{杉本 典充}

Debianの昨今の話題についてのQuizです。

今回の出題範囲は\url{debian-devel-announce@lists.debian.org} や \url{debian-news@lists.debian.org}に投稿された
内容などからです。

\begin{multicols}{2}
%; whizzy-master ../debianmeetingresume201311.tex
% $B0J>e$N@_Dj$r$7$F$$$k$?$a!"$3$N%U%!%$%k$G(B M-x whizzytex $B$9$k$H!"(Bwhizzytex$B$,MxMQ$G$-$^$9!#(B
%

\santaku
{LTS$B%5%]!<%H$,;O$^$C$?(BDebian 7 (wheezy)$B!#(BDebian 6 (squeeze)$B$+$i%5%]!<%H$9$k%"!<%-%F%/%A%c!<$,A}$($^$7$?!#A}$($?%"!<%-%F%/%A%c$N@5$7$$AH$_9g$o$;$O$I$l$G$7$g$&$+!)(B}
{armel $B$H(B armhf}
{armel $B$H(B arm64}
{mipsel $B$H(B ppc64el}
{A}
{wheezy$B$N(BLTS$B%5%]!<%H3+;OD>A0$K(Barmel/armhf$B$r(BLTS$B%5%]!<%H$KDI2C$9$k$+5DO@$,?J$_!"$a$G$?$/(Barmel/armhf$B$,(BLTS$B%5%]!<%H$5$l$k$3$H$K$J$j$^$7$?!#5DO@$O0J2<$N%9%l%C%I$X$I$&$>!#(B \url{https://lists.debian.org/debian-lts/2016/04/msg00045.html}}

\santaku
{Debian GNU/Linux$B$G(Bapt$B$K$h$kG[I[$,;O$^$C$?(Bzfsutils-linux$B%Q%C%1!<%8!#$I$N%3%s%]!<%M%s%H$GDs6!$5$l$F$$$k$G$7$g$&$+!)(B}
{a main}
{b contrib}
{c non-free}
{B}
{ZFS$B$N%=!<%9%3!<%I$O(BCDDL$B$H$$$&%i%$%;%s%9$G8x3+$5$l$F$$$k$?$a!"(BGPLv2$B$G$"$k(BLinux$B%+!<%M%k$N%=!<%9%3!<%I$K%^!<%8$7$F(B(linux$B$H$7$F(B)$B:FG[I[$9$k$3$H$,$G$-$J$$!"$H$5$l$F$-$?7P0^$,$"$j$^$9!#$D$$$K%i%$%;%s%9LdBj$,2r7h$5$l!"(Bcontrib$B$H$7$F(BZFS$B$N%=!<%9%3!<%I$rG[I[$7(Bdkms$B$K$F%S%k%I$7$FMxMQ$9$k7A$KMn$ACe$-$^$7$?!#>\:Y$O0J2<$N(Bweb$B%Z!<%8$H%Z!<%8Fb$N%j%s%/$r$I$&$>!#(B \url{https://bits.debian.org/2016/05/what-does-it-mean-that-zfs-is-in-debian.html}}

\end{multicols}


\dancersection{最近のDebian関連のミーティング報告}{杉本 典充}

\subsection{第138回東京エリアDebian勉強会}

2016年3月5日(土)に第138回東京エリアDebian勉強会を開催しました。会場は日本橋にあるサイボウズさんをお借りして行いました。参加者は14名でした。

発表者はkenhysさん、wskokaさん、Adrianさんの3名でDebian Portsに関連する内容でした。

kenhysさんによる「Debianのインフラを借りるには」では、Debian Projectで用意しているPorterboxという様々なCPUを搭載したサーバ群(または機器群)へログインしてバグ修正や開発を行うために必要な手続きと注意点を発表しました。

wskokaさんによる「tilegx について」では、昨年から実施していたtilegxというDebianにはまだportされていないCPUでDebian GNU/Linuxを動かそうとする試みを発表しました。tilegxアーキテクチャなCPUを搭載したルータの実機上で開発を進めており、72コアのマシンで並列ビルド処理をするデモが行われました。

Debian DeveloperであるAdrianさんによる「Introduction to Debian Ports」では、今回の勉強会のために大阪より駆けつけていただきました。発表は英語で行われ、なぜそのCPUを利用したマシンでDebianを動かそうとするのかというモチベーションや、DebianにおけるPortsを利用する仕組みや新たにPortsを創るときに流れを発表しました。

発表終了後は、各自の課題をこなす作業時間としました。

また勉強会終了後に親睦を深めるため、会食を行いました。

% % (query-replace-regexp "<.*?>" "")
% % (query-replace-regexp "^[	 ]\+" "")

%-------------------------------------------------------------------------------
\dancersection{Buffalo Linkstation 向け Debian Installer}{Roger Shimizu}
%-------------------------------------------------------------------------------

\subsection{始めに}
VGA/HDMI port や Keyboard など付いてない機器、例えばサーバ向けてPC、それから ARM 開発ボードなどの機器に、どうやって Debian を入れておくでしょうか？
RAID 付きの Buffalo Linkstation NAS を例として、紹介致します。

\subsection{Buffalo Linkstation の歴史}

\begin{itemize}
\item 第０世代: Linkstation / Kuro-Box
\item 第１世代: Linkstation HG / Kuro-Box HG
	\begin{itemize}
	\item PowerPC architecture
	\item IDE/PATA のみ (SATA がなし)
	\item 現在あっても使いづらいと思います
	\end{itemize}
\item 第２世代\footnote{MIPSのModelも出たけれど、Model数が少ないし、出てすぐデスコンになってしまったので、こちら省略させて頂きます。}: 
Kuro-Box Pro / Linkstation Live / Linkstation LS-GL/LS-WTGL/LS-WSGL/LS-QL
	\begin{itemize}
	\item ARM architecture
		\begin{itemize}
		\item Debian Etchまで: arm OABI (Old ABI)
		\item Debian Lennyから: armel EABI (new Embedded ABI)
		\end{itemize}
	\item Marvell orion5x 5182 chipset; SATA interface
	\end{itemize}
\item 第３世代: Linkstation LS-XHL/LS-CHL/LS-WXL \& Linkstation LS-VL/LS-WVL/LS-QVL, etc
	\begin{itemize}
	\item Marvell kirkwood 6281 / 6282 chipset
	\item armel architecture なので、第２世代と rootfs の互換性があり、また Debian Kernel 4.4 から同じ「-marvell」との flavour で共通に対応されます。
\footnote{HDDを異なる型番の機器に入れる前に、flash-kernelで適切なDTBをuImage.buffaloに入れ置かないと行けません。}
	\end{itemize}
\item 第４世代: LS-210/LS-220/LS-410/LS-420, etc
	\begin{itemize}
	\item Marvell armada-370 chipset
	\item armhf architecture (hard-float)
	\item 残念ながら、Linux Kernel の方がまだ対応されてないようです。
	\end{itemize}
\end{itemize}

\subsection{Debian Installer の紹介}

\begin{itemize}
\item Debian Installer (略は D-I となります)では、様々機器に Debian をインストールしてくれるツールとなります。
	\begin{itemize}
	\item 非 Linux な OS でも対応される、例え kFreeBSD や GNU/Hurd など
	\item メディアは CD/DVD に限られず、PXE netboot や u-boot など様々柔軟なインストールメディアを提供されております。
	\item モニターなど表示デバイスが付いてない機器でも対応される──network-console イメージ\\
	　⇒ ちょうど Buffalo Linkstation など NAS 機器に向け
	\end{itemize}
\item 現在 D-I に対応されている Linkstation リスト\footnote{d-i daily image が含まれる}:
	\begin{itemize}
	\item Kuro-Box Pro / Linkstation Pro/Live
	\item Linkstation LS-GL / LS-WSGL / LS-WTGL
	\item Linkstation LS-XHL / LS-CHLv2 / LS-WXL / LS-WSXL / LS-VL / LS-WVL / LS-QVL\\
	　⇒ armel の第２世代と第３世代はほとんど対応されてます
	\end{itemize}
\end{itemize}

\subsection{Linkstationへ Debian の インストール仕方}

\subsubsection{Linkstation に Debian Installer を起動させる}
Buffalo Linkstation では、u-boot という boot loader が、1番目の partition (/dev/sda1 又は /dev/md0) に保存される kernel と initrd を読み込んで起動を行います。
それから、D-I イメージを1番目の partition に置いとけば、Debian Installer が起動されます。
手順は下記の通りとなります。
\begin{itemize}
\item 1番目の partition をフォーマットする\footnote{1番目の partition は ext2/ext3 に限られます。}(既に存在するなら省略可)
\item D-I kernel/initrd イメージを、1番目の partition にコピーする。D-I イメージは下記URLでダウンロード出来ます。
	\begin{itemize}
	\item https://d-i.debian.org/daily-images/armel/daily\\/orion5x/network-console/buffalo
	\item https://d-i.debian.org/daily-images/armel/daily\\/kirkwood/network-console/buffalo
	\end{itemize}
\item D-I kernel/initrd image を入れた HDD を Linkstation に取り付けます。
\item Linkstation を起動し、DHCP で IP Address を割り当てられるまで暫く待って、それから割り当てられた IP Address に SSH で接続します。
\item 画面に沿って、通常な Debian Install 行うことが出来ます。
\end{itemize}

\subsubsection{詳細な手順: HDD Partition の作成}
LinkstationのHDDをLinkstationから外し、SATA-USBアダプターなど方法経由でPCに接続します。例えば、その HDD が /dev/sdc として説明します。
LS-WXL/WSXL/WVL/QVL など RAID 機器の場合は、すべての HDD を同じ partition にしても良いです。
\begin{commandline}
$ sudo parted /dev/sdc
(parted) mklabel gpt
(parted) mkpart boot 2048s 1024MiB
(parted) mkpart root 1024MiB 6144MiB
(parted) mkpart swap 6144MiB 6400MiB
(parted) mkpart data 6400MiB -1
# 下記のコマンドは RAID 構成の機種だけ必要となります。
(parted) set 1 raid on
(parted) set 2 raid on
(parted) set 3 raid on
(parted) set 4 raid on
\end{commandline}

\subsubsection{詳細な手順: HDD Partition の確認}
\begin{itemize}
\item Partition を作ったら、確認するとこうなります。
\end{itemize}
\begin{commandline}
(parted) print
Model: SAMSUNG HM250HI (scsi)
Disk /dev/sdc: 250GB
Sector size (logical/physical): 512B/512B
Partition Table: gpt
Disk Flags:
Num Start   End     Size    File sys  Name  Flags
 1  1049kB  1074MB  1073MB            boot  raid
 2  1074MB  6442MB  5369MB            root  raid
 3  6442MB  6711MB  268MB             swap  raid
 4  6711MB  250GB   243GB             data  raid

# 最後にpartedを終了させます。
(parted) quit
\end{commandline}
\begin{itemize}
\item RAID 構成の場合は他の HDD も同じようにセットしてください。
\end{itemize}

\subsubsection{詳細な手順: boot image のセット}
\begin{itemize}
\item boot image kernel/initrd を第１番目 partition にコピーします。(例え、LS-WXL のイメージ)
\item RAID の複数HDDの場合は、念のためすべてHDDをコピーしましょう。\footnote{この時点に「mdadm --create」で RAID の構成をしなくても良いです。Installの際に RAID の作成を行います。}
\end{itemize}
\begin{commandline}
$ sudo mkfs.ext3 /dev/sdc1
$ sudo mount /dev/sdc1 /mnt
$ wget https://d-i.debian.org/daily-images/armel\
  /daily/kirkwood/network-console/buffalo/ls-wxl\
  /uImage.buffalo
$ wget https://d-i.debian.org/daily-images/armel\
  /daily/kirkwood/network-console/buffalo/ls-wxl\
  /initrd.buffalo
$ sudo cp *.buffalo /mnt
$ sudo umount /mnt
\end{commandline}

\subsubsection{詳細な手順: SSH で接続}
\begin{itemize}
\item HDDs を Linkstation に戻して、起動させます。
\item 暫く待つと、Android/iOS アプリ「Fing」のような IP/port scanner で Linkstation に割り当てた IP を見つけます。また、DHCP サーバ側のログでも
\item IP を分かると、SSH コマンドを叩くと debian installer 画面が出て来る：
\end{itemize}
\begin{commandline}
$ ssh installer@<IP address of Linkstation>
\end{commandline}
\begin{itemize}
\item D-I のデフォルトパスワードは「install」となります。
\item command line で操作や log 確認などのため、もう一本 SSH を接続しても良いです。
\end{itemize}

\subsubsection{RAID 構成向け Install 時の注意事項}
\begin{itemize}
\item LS-GL/CHL/XHL/VL などRAID構成ではない機種だとスキップください。
\item もし RAID の設定が見つからない場合は、D-I Partman の画面から一旦「バック」し、「Download installer components」に partman-md や sata-modules などモジュールを選択してから、出来るようになります。
\item もし D-I で RAID を新規に作成する場合、一番目の /dev/md0 は metadata=0 (version 0.90) に設定しないと再起動しなくなります。原因は u-boot が第１番目 partition の kernel/initrd を読み込まなくなるため。現在 partman-md に設定出来なくて\footnote{\url{https://bugs.debian.org/815569}}、command line にしましょう：
\begin{commandline}
# mdadm --create /dev/md0 --level=1 --raid-devices=2 \
  --metadata=0 /dev/sda1 /dev/sdb1
\end{commandline}
\item または (他の HDD が後にします)
\begin{commandline}
# mdadm --create /dev/md0 --level=1 --raid-devices=2 \
  --metadata=0 /dev/sda1 missing
\end{commandline}
\item /dev/md0 以外の RAID は partman-md で設定しても良いです。
\item RAID を新規に作成する (create) と同期化の作業がすぐに始めされ、とても重くて、進行中 Debian Install の作業に影響ならないように /dev/md0 以外の同期化の速度を制限かけた方が良いです：
\end{itemize}
\begin{commandline}
# echo 100 > /sys/block/md{1,2,3}/md/sync_speed_max
\end{commandline}
\begin{itemize}
\item インストールが終わって再起動したら、RAID の同期化の作業は自動的に再開されるので、その後の作業は特に必要ありません。
\end{itemize}

\subsubsection{Debian Install 後の設定}
\begin{itemize}
\item u-boot 環境変数を確認・変更のため、コマンド fw\_printenv / fw\_setenv の設定
\footnote{u-boot 環境変数を変更すると起動しない場合があり、修復手段があまりなくて、気をつけましょう。}
	\begin{itemize}
	\item ほとんどの機種は下記の設定が済みです：
	\end{itemize}
\begin{commandline}
$ sudo echo /dev/mtd2 0x00000 0x10000 0x10000 \
	> /etc/fw_env.config
\end{commandline}
	\begin{itemize}
	\item Kuro-Box Pro は mtd/flash の数が異なるので、下記の設定にしてください。
	\end{itemize}
\begin{commandline}
$ sudo echo /dev/mtd5 0x00000 0x10000 0x10000 \
	> /etc/fw_env.config
\end{commandline}
\item 他の機器で Linkstation の起動ログを確認するため(あるいは、再起動不能の際にデバッグ手段として)、netconsole の設定
	\begin{itemize}
	\item Linkstation 側の設定：
	\end{itemize}
\begin{commandline}
$ sudo cat << EOT >> /etc/initramfs-tools/modules
marvell
mv643xx_eth
netconsole netconsole=@192.168.11.5/,6666@192.168.11.1/
mvmdio
EOT

$ sudo update-initramfs -u
\end{commandline}
	\begin{itemize}
	\item 他の機器でログ収集するコマンド：
	\end{itemize}
\begin{commandline}
$ sudo ip a add 192.168.11.1/24 dev eth0
$ nc -l -u -p 6666 | tee ~/netconsole.log
\end{commandline}
\end{itemize}


\subsection{終わりに}

Kernel の Device-Tree を対応したり、Debian-Installer にパッチを投げたり、Linkstation に Debian Install はやっと出来るようになりました。
今後は Debian Installer を引き続き改善・進化を行って行きたいと思います。
\begin{itemize}
\item Debian Installer に GNU/screen 又は tmux を対応する
\footnote{\url{https://lists.debian.org/debian-boot/2016/02/msg00547.html}}\footnote{\url{https://bugs.debian.org/819397}}\footnote{\url{https://bugs.debian.org/819988}}
	\begin{itemize}
	\item SSH 接続が切れても、installer がバックグラウンドに回せて、SSH で再接続すると元の状態に再開できます。
	\item シェルやログなどより便利にアクセス出来ます。
	\end{itemize}
\item partman-md へ RAID の metadata 指定できるように、対応する\footnote{\url{https://bugs.debian.org/815569}}。
\item 第４世代、armhf/armada-370 の Linkstation を対応する
\end{itemize}


\cleartooddpage

\vspace*{15cm}
\hrule
\vspace{2mm}
\includegraphics[width=2cm]{image200502/openlogo-nd.eps}
\noindent \Large \bf Debian 勉強会資料\\
\noindent \normalfont \debmtgyear{}年\debmtgmonth{}月\debmtgdate{}日 \hspace{5mm}  初版第1刷発行\\
\noindent \normalfont 東京エリア Debian 勉強会 （編集・印刷・発行）\\
\hrule

\end{document}
